\section{Introduction}

Health-regulatory entities are linked, but a natural graph representation does not by itself show that graph learning improves prediction. Nursing homes may share owners, prescribers may share specialties, and products may share reported problem categories. Comparisons must therefore make clear both the information each method uses and the outcome it predicts.

HealthGraphBench addresses these questions with public health-regulatory data\cite{hgb2026}. Version 0.2 preserves the MAUDE and CMS experiments from v0.1 and adds Medicare Part D. We also examine how MAUDE results vary with training duration and with removal of neighbor aggregation.

The central question is when information from related entities adds predictive value, and whether learned message passing improves on simpler baselines. These predictive comparisons are distinct from clinical or operational utility, which this study does not measure.

The contribution is a versioned resource of three task definitions, evaluation protocols, and results, accompanied by a targeted review of methodological context. It supports task-specific conclusions, not a general estimate of graph performance.

\section{Related work: from performance to comparison conditions}

\subsection{Scope of the literature review}

Our targeted literature review through September 2026 covers relational benchmarks, temporal evaluation, simple baselines, and machine-learning study validity. It prioritizes original publications and official documentation; it provides context rather than an exhaustive estimate of how often graph methods succeed.

\subsection{General benchmarks and relational data}

Open Graph Benchmark organized evaluation around shared datasets, splits, and metrics\cite{hu2020ogb}. RelBench shifts attention to prediction on relational databases, where multiple tables must be used jointly\cite{robinson2024relbench}. RelBench v2 extends this approach and introduces, among other things, completion tasks with temporal constraints\cite{gu2026relbench}. These works report benefits from relational learning under their own experimental conditions; their results underscore that performance depends on the task and evaluation design.

HealthGraphBench focuses on precisely defined regulatory tasks, interpretable relational baselines, and explicit limits on inference, complementing the broader benchmarks discussed above.

\subsection{Time, candidates, and simple baselines}

Temporal Graph Benchmark shows that performance varies substantially across datasets and that some simple methods remain competitive\cite{huang2023tgb}. Poursafaei and colleagues also emphasize the role of evaluation strategies and propose a link-memorization baseline\cite{poursafaei2022}. These lessons motivate examining the candidate universe before interpreting a score.

Because MAUDE and Part D exclude links already observed for the target entity, link memorization alone is an inadequate baseline. Peer activity and category popularity may still be predictive. Ranking all eligible candidates avoids dependence on a small sample of easy negatives, but an unrecorded link is not necessarily a true negative.

The fair-comparison requirements studied by Errica and colleagues also remind us that architecture cannot be separated from protocol and reference methods\cite{errica2020}. Their experiments concern graph classification, not the present tasks. We draw a methodological principle from their work, not a direct numerical comparison.

\subsection{Scientific validity and transparency}

Kapoor and Narayanan describe forms of information leakage that can compromise results\cite{kapoor2023leakage}; REFORMS proposes a broader framework for machine-learning studies\cite{kapoor2024reforms}. We use this work to clarify targets, units, splits, comparators, and limits of inference, without claiming full checklist compliance.

Correct chronology is necessary but insufficient. It guarantees neither the historical public availability of every field nor that choices were not adapted after results were examined. Both limits are especially important for registries whose observations may be published or corrected after the event they describe.

\section{Methods}
\label{sec:methods}

\subsection{Study design and data sources}

The study combines the retained MAUDE and CMS experiments from v0.1 with the Part D extension from v0.2.

The duration analysis fits GraphSAGE independently at 0, 3, 10, and 30 epochs on the same 2023 validation cohort. Among the trained candidates, 30 epochs was selected by micro R@10, and only this duration was evaluated in the duration-analysis test runs.

A follow-up compares this selected mean-aggregation model, reused without retraining, with a separately trained self-only model that omits neighbor aggregation. Both use the same 30-epoch duration, fixed before the self-only test runs, but their active model capacities differ. Supplement S11 gives the configurations and provenance.

The new analyses use validation and test periods examined earlier in development, so they are exploratory rather than independent confirmation, even though the model choices were locked on validation before their new test fits. Search budgets across model families are not matched.

\subsection{Tasks, populations, and information availability}
\label{sec:tasks}

\paragraph{MAUDE.} For a product code, predict relationships with problem codes first observed within the retained window. The clock is \texttt{DATE\_RECEIVED}, not a verified publication date for every field. A product must have at least one earlier report; a problem must already exist globally; and the relationship must be absent from the product's history. Thus “threshold 1” means prior product-report support, not severity. All eligible problems are ranked. Novel problems outside the prior vocabulary are outside this ranking task\cite{hgbcontract,hgbcode}.

\paragraph{CMS nursing homes.} For retained Health Standard inspections of US nursing homes, predict whether a standard G--L deficiency is observed at the target inspection using information available before it. The unit is a (CCN, date) episode with at least one earlier inspection. The score is reconstructed retrospectively just before that episode: the task is conditional on the inspection occurring. It predicts neither inspection occurrence nor timing and is not a start-of-year ranking of all facilities. Current-provider extracts and the small number of retained cycles create selection and historical truncation\cite{cmsnursing,hgbcontract,hgbcode}.

\paragraph{Medicare Part D.} Drug identity is the exact generic name after trimming surrounding spaces; it is not normalized to an RxNorm concept. Each target cohort contains at most 2,000 previously observed prescribers, selected by (SHA-256 of NPI, NPI). All drugs in the cohort's prior history are ranked, except those already associated with the target prescriber. A positive is a newly \emph{published and observed} link, not a first prescription. The contract notes CMS suppression of combinations with at most ten claims: an absence is not a clinical negative\cite{cmspartd,hgbpartd}.

\begin{table}[!htbp]
\centering\footnotesize
\caption{Recoverable counts for evaluated splits. MAUDE observations are product--quarters with a positive; CMS observations are retained target inspections; Part D observations are eligible prescriber--years. Units used for interval calculations are not national sample sizes. Sources: reports, denominator audit, annual analysis, and descriptor extraction\cite{hgbresults,hgbanalysis,hgbpartd}; details in S3.}
\label{tab:population}
\begin{tabularx}{\linewidth}{@{}lrrX@{}}
\toprule
Task / split & Observations & Positives & Documented distinct units\\
\midrule
MAUDE, validation 2023 & 3,483 & 7,705 links & 1,641 products\\
MAUDE, test 2024--2025 & 6,370 & 13,174 links & 2,026 products, clusters\\
CMS, validation 2023 & 2,661 & 408 & 2,660 CCNs\\
CMS, test 2024--2025 & 18,985 & 2,423 & 13,889 CCNs, clusters\\
Part D, validation 2023 & 2,000 & 5,794 links & 1,035 with positives\\
Part D, test 2024 & 2,000 & 5,476 links & 1,019 with positives\\
\bottomrule
\end{tabularx}
\end{table}

Supplement S3 reports Part D history, training, and candidate counts. For CMS, constant prevalence scores and the expanding-window protocol reconstruct 2,066, 4,727, and 13,636 training rows before 2023, 2024, and 2025, including 244, 652, and 1,905 positives; this reconstruction is not a direct reading of rows.

At the three historical refits, tabular training sets contained 96,906--124,322 rows, and the graphs 67,676--81,891 edges. Frozen loops imply 270,704--327,564 BPR and 203,028--245,673 GraphSAGE steps; these historical counts are reconstructed from code, not optimizer logs.

For the historical models, reconstructed inventories cover 97.56\,\% of test candidate pairs and 97.80\,\% of positive edges, but only 1,414 of 6,370 positive-containing product--quarters (22.20\,\%) have both nodes represented for every candidate. Supplement S10 reports quarterly counts. This historical inventory analysis measures node-key presence; it does not establish how zero-score fallback affected historical performance. The new models and inspected checkpoints are documented separately in S11.

\begin{table}[!htbp]
\centering\footnotesize
\caption{Evaluation timing and units. In all three tasks, event chronology does not guarantee that every attribute was historically publicly available.}
\label{tab:time}
\begin{tabularx}{\linewidth}{@{}lXXX@{}}
\toprule
 & MAUDE & CMS nursing homes & Part D\\
\midrule
Initial history & 2019Q1--2022Q4 & 2019--2022 & 2019--2022\\
Validation / test & 2023 / 2024--2025 & 2023 / 2024--2025 & 2023 / 2024\\
Scoring time & Before target quarter, by report clock & Just before retained target inspection & Before target service year, reconstructed retrospectively\\
Updates & After each scored quarter; annual refits & Training on strictly earlier years & 2019--2023 history for the 2024 test\\
Horizon / target & Quarterly links & Deficiency at that episode, without validated deployment lead time & Links published in service year\\
\bottomrule
\end{tabularx}
\end{table}

MAUDE is a \emph{rolling} evaluation: once scored, a test quarter enriches the history used subsequently. Refits occur at 2023Q1, 2024Q1, and 2025Q1 using earlier data. Neither model nor history remains immutable throughout 2024--2025. Protocol-defined updating differs from adapting scientific choices after examining scores. For CMS, an ownership-association start date before inspection does not prove that the association was publicly known then. For Part D, service year is not file-publication time\cite{hgbcode,hgbresults}.

\subsection{Heuristics and models actually compared}
\label{sec:models}

\paragraph{MAUDE: neighbor frequency.} Let $H_p$ be the set of problems previously associated with product $p$, and $U_d$ the set of products historically associated with problem $d$. Neighbors are $N_p=(\bigcup_{d\in H_p}U_d)\setminus\{p\}$. For candidate $c\notin H_p$, the two scores are
\begin{equation}
 s_{\mathrm{global}}(p,c)=|U_c|,\qquad
 s_{\mathrm{neighbor}}(p,c)=|N_p\cap U_c|.
\end{equation}
Each distinct peer counts once: neither report-volume weighting nor degree normalization is applied. Missing support gives zero. Ties are resolved by descending global support, then ascending problem code\cite{hgbcode}. This definition clarifies what the heuristic personalizes.

The retained historical MAUDE comparators are logistic regression, boosted stumps, rank-8 spectral factorization, BPR with fixed one-hop averaging, and a \emph{one-hop GraphSAGE variant}\cite{hamilton2017,rendle2009,hgbcode}. The historical variant uses trainable node inputs, two shared transforms, mean neighbor aggregation, a $\tanh$ activation, and a dot-product score; its configuration is dimension 8, three epochs, eight neighbors, learning rate 0.02, and regularization 0.0005. The additional training-duration and aggregation-free comparisons are distinguished from this retained configuration in the Results and Supplement.

\paragraph{Part D: specialty and overlap.} Specialty popularity counts historical prescribers with the same specialty who are associated with the drug. Specialty is the unique nonblank value in the latest observed year; otherwise, scoring falls back to global support. For prescriber histories $D_i,D_j$, overlap uses
\begin{equation}
 s_{\mathrm{overlap}}(i,d)=\sum_{j\ne i}\frac{|D_i\cap D_j|}{|D_i\cup D_j|}\,\mathbf{1}\{d\in D_j\}.
\end{equation}
Only peers with nonempty overlap contribute; the sum is not normalized by the number of peers. Zero-score ties follow the shared rule: descending global support, then ascending canonical key. The learned models are logistic regression on ten historical features and a prescriber--drug BPR augmented with a specialty representation. This BPR does not propagate messages\cite{hgbcode,hgbpartd}.

\paragraph{CMS: owner augmentation.} Both models are logistic regressions: eleven local numeric features and state indicators, then twelve additional relational aggregates. The \emph{combined} representation prioritizes CHOW owner identifiers; normalized current names are a fallback for CCNs absent from that representation. The two sets are not systematically combined. Peers are other CCNs sharing an association that starts on or before the target date. Since end dates are not all known, this rule does not establish legally active ownership. Peer inspection and deficiency counts are accumulated strictly before the target; rates divide total deficiencies by total inspections, with zero when the denominator is empty\cite{hgbcode}.

CMS models use the same annual fitting procedure, with 300 iterations, learning rate 0.08, and $L_2=0.05$, standardized on training data. These are retained settings, not evidence of optimality; comparable search budgets and per-model resource trade-offs have not been established.

\subsection{Metrics, ties, and uncertainty}
\label{sec:metrics}

At cutoff $K$, we denote micro recall by $R@K$. For observation $i$ with $P_i>0$ positives and $h_i(K)$ recovered links, micro recall is $\sum_i h_i(K)/\sum_i P_i$ and macro recall is the mean of $h_i(K)/P_i$. MRR averages the reciprocal rank of the first positive. For Part D, precision and burden also include zero-positive observations: with $a_i(K)=\min(K,|C_i|)$, precision is $\sum_i h_i(K)/\sum_i a_i(K)$ and mean burden is $\sum_i a_i(K)/n$. A ranking cutoff is not calibrated confidence.

CMS reports ROC AUC, Brier score, top-decile precision and recall, and two distinct average-precision conventions. \emph{Historically tie-broken rank AP}, denoted $\mathrm{AP}_{\mathrm{rank}}$, averages precision at positive ranks after sorting by decreasing score and increasing row index (date then CCN). \emph{Score-threshold AP}, denoted $\mathrm{AP}_{\mathrm{thresholds}}$, groups exactly equal scores. If threshold group $j$ contains $t_j$ positives, with cumulative $T_j$ positives among $N_j$ observations,
\begin{equation}
 \mathrm{AP}_{\mathrm{thresholds}}=\sum_j\frac{t_j}{P}\frac{T_j}{N_j}.
\end{equation}
This is non-interpolated precision averaged over thresholds\cite{sklearnap}. AUC uses average ranks for ties. The historical decile retains exactly $\lceil0.1n\rceil$ rows. Historical AP values and confidence intervals refer only to $\mathrm{AP}_{\mathrm{rank}}$; no interval is transferred to $\mathrm{AP}_{\mathrm{thresholds}}$.

The historical v0.1 intervals and two complementary post-hoc paired bootstraps apply to retained predictions, not to the new MAUDE analyses. Each bootstrap resamples products (MAUDE) or prescribers (Part D), retaining all observations of each sampled entity over 1,000 draws. The intervals are conditional on those predictions and populations; they do not cover selection, model fitting, or all dependence between entities, and are not independent validation.

\section{Results}

\subsection{Historical MAUDE baseline: neighbor frequency at ten, cutoff dependent}

\begin{table}[!htbp]
\centering\small
\caption{Retained historical v0.1 MAUDE test results, 2024Q1--2025Q4, prior product-report support $\geq1$. R@K is micro recall. Bold marks the highest historical point value in a column; it is not an additional superiority test. Source: unified table\cite{hgbresults}.}
\label{tab:maude}
\input{tables/maude_en.tex}
\end{table}

In the historical results, neighbor frequency has R@10 of 0.192045, compared with 0.184454 for global popularity and 0.172840 for the three-epoch one-hop GraphSAGE variant (Table~\ref{tab:maude}). The post-hoc product bootstrap for neighbors minus global gives $+0.007591$ (95\,\% CI $[+0.005497,+0.009820]$). The published GraphSAGE interval also applies only to the historical predictions; neither interval describes the new analyses. Macro recall and MRR show the same historical ordering\cite{hgbresults}.

At $K=20$, spectral factorization reaches 0.293912 versus 0.293077 for neighbor frequency, a descriptive difference of 0.000835, with no interval (Figure~\ref{fig:cutoffs}). Thus the leading method depends on the selection budget.

\begin{figure}[!ht]
\centering
\includegraphics[width=\linewidth]{figures/maude_cutoffs_en.pdf}
\caption{Historical MAUDE recall differences relative to neighbors. The R@10 intervals are the retained GraphSAGE--neighbor interval and post-hoc global--neighbor interval; the latter is plotted with the sign reversed from the neighbors-minus-global estimate in Table~\ref{tab:ci}. The K=20 point values compare spectral factorization and neighbor frequency.}
\label{fig:cutoffs}
\end{figure}

\subsection{Exploratory MAUDE training-duration and aggregation comparisons}
\label{sec:maude-post-rc1}

The duration analysis selected 30 epochs on 2023 validation, yielding micro R@10 of 0.205062. The selected model then achieved pooled 2024--2025 R@10 of 0.210642, recovering 2,775 of 13,174 positive links. This point estimate exceeds historical neighbor frequency (0.192045) and three-epoch GraphSAGE (0.172840). However, the new test runs evaluated only the selected duration: the comparison with historical three-epoch GraphSAGE is not a matched duration experiment.

The aggregation comparison reused this mean model and trained a self-only variant for the same cohorts and duration. Pooled test R@10 was 0.210642 for mean aggregation and 0.035600 for self-only, a mean-minus-self-only difference of $+0.175042$. Supplement S11 reports validation and annual results, denominators, and optimization traces.

The self-only variant has no message passing and is not a GNN. It uses 64 active shared-transformation coefficients, compared with 128 for mean aggregation, alongside the same number of trainable node-vector coefficients. Its loss stayed close to $\ln 2$ in all three fits. Thus the observed difference is specific to these configurations: it neither isolates the effect of aggregation nor establishes that the self-only architecture cannot learn under other settings. The evaluated cohorts contain only product--quarters with observed positives.

\subsection{CMS inspections: an uncertain increment from owner aggregates}

\begin{table}[!htbp]
\centering\small
\caption{CMS inspections, pooled 2024--2025 test predictions. $P_{10\%}$ and $R_{10\%}$ denote precision and recall in the top decile ranked across pooled scores, not separately by year. Brier should decrease; other metrics should increase. Source: unified table\cite{hgbresults}.}
\label{tab:cms}
\input{tables/cms_en.tex}
\end{table}

Local history reaches AUC 0.619975 and owner augmentation 0.621293 (Table~\ref{tab:cms}). The difference of 0.001318 has a historical interval containing zero. AP rises only slightly, while Brier and pooled top-decile precision deteriorate slightly in point value. The augmentation therefore does not provide a consistent benefit across the reported measures\cite{hgbresults}.

No equivalence margin was specified, so the small difference and interval crossing zero do not establish equivalence. Annual AUC differences were $+0.000545$ in 2024 and $-0.000087$ in 2025; pooled AUC is not the simple average of annual AUCs\cite{hgbresults}.

\subsection{CMS: annual discrimination and pooling}
\label{sec:cmsannual}

\begin{table}[!htbp]
\centering\footnotesize
\caption{Annual CMS test results for 2024 and 2025. Prevalence is 14.0644\,\% and 11.6118\,\%, respectively; AP follows the historical rank convention. The supplement includes 2023 validation\cite{hgbanalysis}.}
\label{tab:cmsannual}
\input{tables/cms_annual_test_en.tex}
\end{table}

The 8,909 inspections in 2024 include 1,253 positives; the 10,076 inspections in 2025 include 1,170. Calculated from unrounded AUCs, the owners-minus-history difference is $+0.000545$ in 2024 and $-0.000087$ in 2025, rounded to six decimal places.

Annual top-decile precision/recall pairs $(P_{10\%},R_{10\%})$ are, in 2024: prevalence (0.158249; 0.112530), local history (0.273850; 0.194733), and owners (0.267116; 0.189944); in 2025: prevalence (0.123016; 0.105983), local history (0.237103; 0.204274), and owners (0.235119; 0.202564). Brier increases slightly in both years.

Annual results better describe within-year discrimination than the pooled result alone.

A descriptive decomposition separates positive--negative pairs from the same year from cross-year pairs. If $P_y,N_y$ are positive and negative counts, $T=(\sum_yP_y)(\sum_yN_y)$, and $W=\sum_yP_yN_y$, then
\begin{equation}
 A_{\mathrm{pooled}}=\frac{W}{T} A_{\mathrm{within}}+
 \frac{T-W}{T} A_{\mathrm{cross}},\qquad
 A_{\mathrm{within}}=\frac{\sum_y P_yN_y A_y}{W}.
\label{eq:decomp}
\end{equation}
Within-year AUC is 0.625299 for local history and 0.625515 with owners, a difference of $+0.000216$. The weighted within- and cross-year contributions to the pooled difference of $+0.001318$ are $+0.000108$ and $+0.001210$. This decomposition is descriptive and carries no interval.

The prevalence reference assigns a constant score within each year, so its annual AUC is 0.5. It assigns a higher score to the 2025 cohort than to 2024, yielding pooled AUC 0.472568; pooled discrimination also reflects cross-year ranking.

\subsection{CMS: the effect of ties on average precision}
\label{sec:ties}
Score-threshold AP differs from historical rank AP for the constant prevalence reference because the latter resolves ties by row order. The threshold values are 0.140644 in 2024, 0.116118 in 2025, and 0.122069 pooled; historical rank AP was 0.148816, 0.118687, and 0.121754, respectively. The difference reflects tie handling rather than added score discrimination.

For local history, score-threshold AP is 0.211180530 in 2024, 0.183722903 in 2025, and 0.196022698 pooled, identical to rank AP. The owner-minus-local differences are $+0.001796527$, $-0.000175565$, and $+0.001249204$, respectively; no interval is available for this metric.

The two learned models have unique decile boundaries in each year and in the pooled test, so boundary tie-breaking does not affect selection. The prevalence reference ties all observations within each year, making row order relevant to decile precision. Under uniform selection, expected precision equals annual prevalence; in the pooled test, the cutoff falls within the 2025 tied group, so the expectation is the 2025 prevalence, as detailed in S7.

\subsection{Part D: cross-domain extension and selection burden}

\begin{table}[!htbp]
\centering\small
\caption{Part D test results for 2024. Recall and MRR cover the 1,019 prescribers with positives; P@10 covers all 2,000 eligible prescribers. Scores use exact generic-name identity\cite{hgbresults,hgbpartd}.}
\label{tab:partd}
\input{tables/partd_en.tex}
\end{table}

At $K=10$, specialty popularity has 0.037071 higher recall than relational BPR; the complementary paired prescriber-level interval for BPR minus specialty is $[-0.047179,-0.027336]$ (Table~\ref{tab:ci}). It also exceeds logistic regression by 0.049671 (Table~\ref{tab:partd}).

Heuristic ordering varies by cutoff: at $K=5$, history overlap reaches 0.128561 versus 0.127465 for specialty. No interval is reported for this contrast. Because there is no strictly local comparator, the result does not isolate a relational-information effect.

Of the 2,000 prescribers, 1,019 had observed positive links. Specialty popularity retrieved 1,193 of 5,476 links in 20,000 suggestions: R@10 was 0.217860 and P@10 was 0.059650. The 981 prescribers without an observed positive also received ten suggestions; these unobserved links are not established clinical errors\cite{hgbpartd}.

\Needspace{10\baselineskip}
\subsection{Uncertainty summary}

The intervals in Table~\ref{tab:ci} apply only to the retained historical contrasts.

\begin{table}[!htbp]
\centering\small
\caption{Historical contrasts and 95\,\% intervals. The first six rows retain v0.1-published intervals; the final two summarize 1,000-draw, post-hoc paired bootstraps on retained predictions. These intervals are conditional on the observed predictions and do not cover selection or training. No interval is available for the exploratory duration and aggregation-free analyses. Sources: v0.1/v0.2 reports and outputs in supplement S6\cite{hgbresults}.}
\label{tab:ci}
\begin{tabularx}{\linewidth}{@{}Xrr@{}}
\toprule
Contrast / metric & Difference & 95\,\% interval\\
\midrule
MAUDE (v0.1): GraphSAGE (3 epochs) -- neighbors, R@10 & $-0.019204$ & $[-0.023578,-0.014845]$\\
MAUDE (v0.1): same contrast, macro R@10 & $-0.027883$ & $[-0.033630,-0.022316]$\\
MAUDE (v0.1): same contrast, MRR & $-0.018949$ & $[-0.023026,-0.014850]$\\
CMS (v0.1): owners -- local, AUC & $+0.001318$ & $[-0.001458,+0.004329]$\\
CMS (v0.1): same contrast, $\mathrm{AP}_{\mathrm{rank}}$ & $+0.001249$ & $[-0.001186,+0.003570]$\\
CMS (v0.1): same contrast, Brier & $+0.000132$ & $[-0.000038,+0.000303]$\\
MAUDE (retained predictions): neighbors -- global, R@10 & $+0.007591$ & $[+0.005497,+0.009820]$\\
Part D (retained predictions): BPR -- specialty, R@10 & $-0.037071$ & $[-0.047179,-0.027336]$\\
\bottomrule
\end{tabularx}
\end{table}

\subsection{Controls and a deferred task}

In the synthetic ownership-topology control, the neighborhood score itself enters the label-generating mechanism and is then evaluated directly on test nodes. Mean AUC was 0.503087 for $\beta=0$ and 0.817456 for $\beta=2.5$ across three seeds\cite{hgbcontrols}. This is a signal-construction diagnostic, not evidence about the CMS regressions, the MAUDE health-data models, or a causal effect; see Supplement S2.4.

A separate synthetic two-block graph used 24 products, 12 problems, 120 training edges, and 24 masked edges, with every target node known. Masked-edge Recall@1 was 0.125000 before training, 0.083333 after three epochs, and 1.000000 after thirty epochs under the same settings apart from duration. The check shows that the extracted implementation can recover a constructed signal; it does not establish that three epochs suffice on MAUDE or identify an aggregation effect. Protocol and gradient checks are in S8.

ClinicalTrials.gov/AACT remains outside the three admitted tasks. Its target is first results publication within 365 days after origin, for interventional trials whose actual primary completion occurred 0--90 days earlier; it is not a legal-compliance label. Sponsor history minus local features gives $\Delta\mathrm{AP}=-0.048773$ in 2022, CI $[-0.088028,-0.021689]$; heterogeneous context minus sponsor history gives $-0.108077$ in 2023, CI $[-0.171780,-0.049169]$\cite{hgb2026,hgbhistory}. These results remain separate by origin and are not pooled across the three tasks.

The ClinicalTrials memo links deferral partly to unstable gains, so task inclusion was not demonstrably independent of prior results. This history tempers confirmatory interpretations and motivates prespecifying task-admission criteria as well as model choices.

\section{Discussion}

\subsection{What the results establish, and what they do not disentangle}

The findings depend on task and comparator. In the historical MAUDE table, neighbor frequency exceeds three-epoch GraphSAGE at R@10; the validation-selected 30-epoch model has a higher point estimate, but its test comparison with the historical model is not matched. The self-only model scores below mean aggregation, yet active capacity differs and its training loss remains nearly flat. CMS shows only a small, uncertain increment from owner aggregates, while Part D specialty popularity outperforms the learned methods tested. These comparisons do not establish a general ranking of graph and non-graph methods.

Several explanations remain possible: relational summaries may capture much of the available signal, message passing may help under some settings, or a chosen representation may fail to use available information effectively. The self-only model's nearly flat loss does not identify why its scores differ. Because its active shared-transform capacity is lower than that of mean aggregation, the observed contrast does not isolate the effect of aggregation.

The CMS comparison directly contrasts owner aggregates added to the same local-history model. For MAUDE and Part D, the methods differ in information, objective, regularization, and capacity; the new self-only comparison narrows the question but does not equalize capacity. Positive results from RelBench remain compatible with these task-specific findings\cite{robinson2024relbench,gu2026relbench}.

\subsection{Three validity limitations}

\paragraph{Internal validity.} The additional MAUDE analyses use periods already examined in development. Existing intervals are conditional on retained predictions and do not quantify variation from training or selection. The fixed configurations and unequal active capacity in the self-only comparison limit attribution; search budgets across model families are not matched.

\paragraph{Measurement validity.} Targets reflect observed reports and records. Underreporting, coding, vocabulary changes, low-volume suppression, and shallow history can alter measured links or deficiencies. Event dates do not establish that all attributes were public at that time; historical publication-time availability for Part D has not been verified\cite{hgbresults}. FDA reporting limitations also preclude interpreting these scores directly as medical risk\cite{fda2026}.

\paragraph{External validity.} Three heterogeneous tasks and a bounded Part D cohort do not represent all public health, care systems, or relational databases. Retained inspection data cover few years; eligibility excludes entities without a history; and the ranking universe excludes objects outside the historical vocabulary. Performance applies to these subpopulations and periods. The benchmark does not measure user decisions based on scores or their consequences.

\subsection{Why bounded, mixed results remain informative}

The resource is most useful as a comparison of task-specific baselines under explicit candidate and time definitions, rather than as a contest with one winner. The MAUDE findings show that the historical ranking should not be generalized across training durations. They also motivate investigating how training configuration and access to neighbors affect learning, without identifying a causal benefit of either. Stronger confirmation would require an unexamined origin, choices documented before evaluation, and evidence that relevant fields were historically available.

Clinical or operational utility and comparable resource trade-offs remain unmeasured.

\section{Using the resource and reproducing analyses}
\label{sec:repro}
HealthGraphBench exposes \texttt{load\_task}, \texttt{Split}, \texttt{PredictionSet}, and the \texttt{fit\_predict} and \texttt{evaluate} methods\cite{hgbinterface}. MAUDE and CMS require external snapshots. The Part D adapter consumes previously computed, verified rankings: calling it is not a new training run. A researcher can inspect a comparator and its metrics under the same contract before changing a method separately.

Each model adapter must preserve its task's periods, candidate construction, denominator, and prediction format; S9 provides examples.

Supplement S9 describes the interface, retained data, and replay procedures; S11 documents the additional MAUDE analyses and their evidence limits. Replay reproduces the retained bootstrap summaries and CMS metrics, but does not rerun the historical models. Full candidate-score streams for the new analyses are not included.


\section{Conclusion}

HealthGraphBench provides temporal task definitions and interpretable baselines for three distinct health-regulatory prediction problems. Its results depend on the comparison and training design: historical three-epoch GraphSAGE trails neighbor frequency, whereas a validation-selected 30-epoch model has a higher point estimate; the self-only control performs poorly under the shared settings but differs in active capacity. CMS owner aggregates yield a small pooled AUC increment with an interval spanning zero, and Part D specialty popularity exceeds BPR. These findings support transparent, task-specific benchmarking and caution against extending a single configuration's ranking to graph methods generally.
